% Bagian: sets-functions-relations
% Bab: relations
% Subbagian: operations

\documentclass[../../../include/open-logic-section]{subfiles}

\begin{document}

\olfileid[id]{sfr}{rel}{ops}
\olsection{Operasi pada Relasi}

Sering kali berguna untuk mengubah atau menggabungkan relasi. Dalam
\olref[sfr][rel][ord]{prop:stricttopartial}, kita meninjau \emph{gabungan}
relasi, yang tidak lain adalah gabungan dua relasi yang dipandang sebagai
himpunan pasangan. Dengan cara serupa, dalam
\olref[sfr][rel][ord]{prop:partialtostrict}, kita meninjau selisih relatif
relasi. Berikut beberapa operasi lain yang dapat kita lakukan pada relasi.

\begin{defn}\ollabel{relationoperations}
Misalkan $R$ dan $S$ merupakan relasi, dan $A$ sembarang himpunan.

\emph{Invers} dari $R$ adalah $R^{-1} = \Setabs{\tuple{y, x}}{\tuple{x,
    y} \in R}$.

\emph{Hasil kali relatif} dari $R$ dan $S$ adalah $(R \mid S) =
\{\tuple{x, z} : \exists y(Rxy \land Syz)\}$.

\emph{Pembatasan} $R$ pada $A$ adalah $\funrestrictionto{R}{A}= R
\cap A^2$.

\emph{Penerapan} $R$ pada $A$ adalah $\funimage{R}{A} = \{y :
(\exists x \in A)Rxy\}$.
\end{defn}

\begin{ex}
Misalkan $S \subseteq \Int^2$ merupakan relasi penerus pada~$\Int$, yakni
$S = \Setabs{\tuple{x, y} \in \Int^2}{x + 1 = y}$, sehingga $Sxy$ jika dan
hanya jika $x + 1 = y$.

$S^{-1}$ merupakan relasi pendahulu pada $\Int$, yakni
$\Setabs{\tuple{x,y}\in\Int^2}{x -1 =y}$.

$S\mid S$ adalah
$ \Setabs{\tuple{x,y}\in\Int^2}{x + 2 =y}$.

$\funrestrictionto{S}{\Nat}$ merupakan relasi penerus pada~$\Nat$.

$\funimage{S}{\{1,2,3\}}$ adalah $\{2, 3, 4\}$.
\end{ex}

\begin{defn}[Penutupan transitif]Misalkan $R \subseteq A^2$ merupakan relasi
biner.

\emph{Penutupan transitif} dari~$R$ adalah $R^+ = \bigcup_{0 < n \in
\Nat} R^n$, dengan $R^1 = R$ dan $R^{n+1} = R^n \mid R$ didefinisikan secara
rekursif.

\emph{Penutupan refleksif-transitif} dari $R$ adalah
$R^* = R^+ \cup \Id{A}$.
\end{defn}

\begin{ex}
Ambil relasi penerus $S \subseteq \Int^2$. $S^2xy$ jika dan hanya jika
$x + 2 = y$, $S^3xy$ jika dan hanya jika $x + 3 = y$, dan seterusnya. Jadi,
$S^+xy$ jika dan hanya jika $x + n = y$ untuk suatu $n \geq 1$. Dengan kata
lain, $S^+xy$ jika dan hanya jika $x < y$, dan $S^*xy$ jika dan hanya jika
$x \le y$.
\end{ex}

\begin{prob}
Tunjukkan bahwa penutupan transitif dari $R$ memang transitif.
\end{prob}

\end{document}
